% Pursuing Stacks (cotes 134-2 to 134-8): archive cote 134-6, site page 3 (PDF page 4), typescript page 322.
% typescript. Draft compared: ps5.tex:1459-1535.
% AI-transcribed by Claude Opus 5.5 (claude-opus-5-5, adaptive thinking, effort high), Message Batches API, batch ps3 (msgbatch_01XdFmuVVBAcgc7EFxdtCqsj), 2026-11-12; not yet checked by eye.
\scanpage{134-6/3}
\mspage{322}

(2) \qquad $D.(Ab) \ = W^{-1}\ Ch.(Ab)$ \qquad ,

where W means ``quasi-isomorphisms'', i.e. maps inducing isomorphisms for
all homology groups. The most common way for defining the canonical
functor (1), where \unsure{as|Hs} usual here (Hot) is defined as $\underline{W}^{-1}$(Cat), is
via the \inserted{test} category \inserted{$\Delta$}, as the composition in the bottom row of
\diagram{
\begin{tikzcd}
(Cat) \ar[r, "\hand{$i^*$}"] \ar[d] & \hand{$\Delta$}^{\wedge} \ar[r, "\hand{$Wh_{\Delta}$}"] \ar[d] & \hand{$\Delta$}^{\wedge}_{ab} \ar[r, "\hand{$\approx\, DP$}"] \ar[d] & Ch.(Ab) \ar[d] \\
(Hot) \ar[r, "\hand{$\approx$}"] & Hot_{\hand{$\Delta$}} \ar[r] & Hotab_{\hand{$\Delta$}} \ar[r, "\hand{$\approx$}"] & D.(Ab)
\end{tikzcd}
} \quad ,
(3)

which is deduced from the top row by passing to the localized categories.
The \xout{\unsure{functors}} subscript ab in $\hand{\Delta}^{\wedge}_{ab}$ \xout{it}denotes the \xout{\unsure{sub}}category of abelian
group objects in $\hand{\Delta}^{\wedge}$ , the functor \xout{Wh}
\[ (4) \qquad Wh_{\hand{\Delta}} : \hand{\Delta}^{\wedge} \to \hand{\Delta}^{\wedge}_{ab} \]
is the ``\uline{abelianization functor}'' obtained by composing a presheaf
$\hand{\Delta}^{o} \to$ (Sets) with the abelianization functor
\[ (Sets) \longrightarrow (Ab) \qquad , \qquad X \longmapsto \underline{Z}^{(X)} \quad . \]
We call this functor \inserted{$Wh_{\Delta}$} also the ``Whitehead functor'', as \over{the}{\unsure{its|it's}} main property
\deleted{of this functor} is expressed in \uline{Whitehead's} \uline{theorem}, namely that it
is compatible with weak equivalences (where weak equivalences in $\hand{\Delta}^{\wedge}_{ab}$
are defined in terms of the underlying semisimplicial sets, forgetting
the addition laws). The localized category of $\hand{\Delta}^{\wedge}_{ab}$ with respect to
the latter notion of weak equivalence is denoted by $Hotab_{\hand{\Delta}}$ , The
functor
\[ \hand{(4')} \qquad Hot_{\hand{\Delta}} \to Hotab_{\hand{\Delta}} \]
induced by $Wh_{\hand{\Delta}}$ may be equally (and still more validly) be designated
by $Wh_{\hand{\Delta}}$. The two subscripts $\hand{\Delta}$ (in Wh and in Hotab) refer to the
fact that the notions make still a sense when $\hand{\Delta}$ is replaced by an ar-%
bitrary small category A\hand{, cf. below.}

The functor DP in the top row is the well-known \uline{Dold-Puppe} \uline{func-}%
\uline{tor}, which is an equivalence of categories. As for $i^*$ , it is de-%
fined in terms of an arbitrary test functor
\[ i : \hand{\Delta} \longrightarrow (Cat) \quad , \]
which may be either the \xout{\illegible{1}}\inserted[typed]{standard} inclusion (which is the more commonly
used one) or the canonical functor $a \to \hand{\Delta}_{/a}$ , called $i_{\hand{\Delta}}$ . The
functors $i^*$ corresponding to different choices of i are of course
in general non isomorphic, however (as follows from \deleted{the end remarks of} \inserted{section 77)}

\margin{3}
